\documentclass[11pt]{article}

\usepackage{amsmath, amsfonts, amssymb, amstext, amscd, amsthm, makeidx, graphicx, hyperref, url, enumitem, geometry, booktabs, multirow}
\usepackage[table]{xcolor}

\newenvironment{question}[1]
    {
        \begin{center}
        \textbf{#1} \rmfamily
        \bigskip
        \hrule
    }
    {
        \end{center}
    }

\setlength{\textheight}{22.5cm}
\geometry{
    lmargin = 1in,
    rmargin = 1in
}
\setlength{\topmargin}{0pt}
\setlength{\headsep}{1pt}
\setlength{\parskip}{1mm}
\setlength{\parindent}{0mm}
\renewcommand{\baselinestretch}{1.03}

\newcommand{\unitfrac}[4]{\frac{#1 \text{ #2}}{#3 \text{ #4}}}


\title{\vspace{-0.7cm}EE/CS 148B - Homework 3}
\date{\today}
\author{Aiden Di Carlo}

\begin{document}

\maketitle

\begin{question}{Vision Transformer}
\begin{enumerate}[label=(\alph*)]

    \item vit\_pooling: For tasks requiring spatial reasoning, attention pooling is expected to perform best because it can learn to weight tokens selectively based on task-relevant spatial cues, rather than treating all positions equally (mean pooling) or relying on a single generic summary (CLS). Mean pooling discards per-patch identity by averaging, which collapses distinctions between, say, a red cube on the left versus the right; the CLS token suffers the same problem more acutely, since it is trained to summarize global semantics and does not preserve fine-grained spatial layout at all. When an entire image is condensed into a single CLS vector before being passed to a language model, all explicit positional relationships between objects are lost. The model receives no direct signal about where each object is, only a holistic blend of what is in the image.

    \item vit\_patch\_size:

    \textbf{(1) Patch counts and attention cost.}
    For a $224\times224$ image the number of patches is $N=(224/P)^2$:

    \begin{center}
    \begin{tabular}{cc}
    \toprule
    $P$ & $N$ \\
    \midrule
    8  & 784 \\
    16 & 196 \\
    32 &  49 \\
    \bottomrule
    \end{tabular}
    \end{center}

    Self-attention scales as $O(N^2 d_\text{model})$, and $N \propto P^{-2}$. Thus, the overall cost scales as $O(P^{-4} d_\text{model})$. Halving the patch size increases attention compute by $2^4 = 16\times$, so shrinking $P$ is extremely expensive.

    \textbf{(2) Wall-clock times} (batch of 16 images, $d_\text{model}=384$, 6 heads, 6 blocks; RTX 4070 Laptop GPU; 5 warmup + 20 timed steps):

    \begin{center}
    \begin{tabular}{ccc}
    \toprule
    $P$ & Mean (ms) & Std (ms) \\
    \midrule
    8  & 84.9 & 1.8 \\
    16 & 14.8 & 1.3 \\
    32 & 10.1 & 1.3 \\
    \bottomrule
    \end{tabular}
    \end{center}

    The empirical ratios ($8.4\times$ between $P=8$ and $P=16$, and $1.5\times$ between $P=16$ and $P=32$) are smaller than the theoretical $16\times$ and $256\times$ because the GPU is compute-bound only for the largest sequence length; at $P=16$ and $P=32$ the bottleneck shifts to memory bandwidth and kernel launch overhead, compressing the observed differences. The $P=8$ overhead, roughly $8\times$ slower than $P=16$, confirms that the quadratic attention term dominates at high patch counts.

    \textbf{(3)} Accept smaller patches when fine-grained spatial fidelity is critical and the compute budget allows it (e.g., OCR, dense prediction, or counting small objects where coarser tokens would merge distinct features).

\end{enumerate}
\end{question}

\begin{question}{CLIP-Style Contrastive Pretraining}
\begin{enumerate}[label=(\alph*)]

    \item infonce: The loss is averaged in both directions because each (image, text) pair should serve as a positive for both retrieval directions---finding the right caption given an image, and the right image given a caption. Using only one direction would bias the model toward one retrieval mode and leave half the pairwise supervision signal in each batch unused.

    \item clip\_train:

    \begin{figure}[h]
    \centering
    \includegraphics[width=0.48\textwidth]{../runs/clip_eurosat/loss_curve.png}
    \includegraphics[width=0.48\textwidth]{../runs/clip_eurosat/acc_curve.png}
    \caption{Left: mean training loss per epoch. Right: zero-shot validation accuracy per epoch. EuroSAT, 20 epochs, A100 GPU.}
    \end{figure}

    Validation accuracy rises steeply during the first eight epochs (53\% $\to$ 87\%), then continues to improve at a diminishing rate before plateauing near 91\% from epoch 17 onward. Training loss, by contrast, decreases monotonically throughout all 20 epochs---including the final epochs where accuracy is essentially flat---indicating the model keeps tightening image--text alignment on training pairs past the point of diminishing returns for zero-shot classification. This is expected under EuroSAT's class-template captions: many examples per batch share the same caption, so the loss can still improve by better separating near-duplicate positives even after the coarser class-level decision boundaries are learned.

    \item clip\_zeroshot:

    \begin{center}
    \includegraphics[width=\textwidth]{../runs/clip_eurosat/zeroshot_examples.png}
    \end{center}

    Nearly all errors are reasonable confusions between visually similar classes: the four agricultural/vegetation categories (Pasture, Herbaceous Vegetation, Annual Crop, Permanent Crop) appear as low-texture green patches at $64\times64$ resolution, making them hard to distinguish, and Industrial Buildings versus Residential Buildings are both regular grids of rooftops separated by roads. The one surprising error is Annual Crop predicted as Highway: at this resolution, the parallel furrows of a ploughed field closely resemble lane markings on a multi-lane road. Because all five wrong predictions still place the correct class in the top-3, the mistakes reflect genuine local ambiguity rather than a failure of the embedding space: the model has learned semantically coherent clusters (``green land,'' ``built environment,'' ``water'') but struggles to separate fine-grained classes within each cluster.

\end{enumerate}
\end{question}

\begin{question}{LoRA Fine-Tuning}
\begin{enumerate}[label=(\alph*)]

    \item lora\_linear: Applying LoRA (rank 8, $\alpha=16$) to every \texttt{q\_proj} and \texttt{v\_proj} in the ViT and adding a 45-way classification head gives the following parameter breakdown:

    \begin{center}
    \begin{tabular}{lr}
    \toprule
    & Count \\
    \midrule
    Total parameters         & 11,013,165 \\
    Trainable parameters     &    275,373 \\
    \quad of which LoRA A/B  &    258,048 \\
    \quad of which head      &     17,325 \\
    Trainable / Total        &      2.50\% \\
    \bottomrule
    \end{tabular}
    \end{center}

    The LoRA count follows from 36 adapted heads (6 heads $\times$ 6 blocks), each contributing $2 \times (r \cdot d_\text{model} + d_\text{head} \cdot r) = 2 \times (8 \cdot 384 + 64 \cdot 8) = 7{,}168$ parameters, giving $36 \times 7{,}168 = 258{,}048$. Total parameters for LoRA exceeds the base model (10,755,117) by exactly this amount.

    \item lora\_compare:

    \begin{center}
    \begin{tabular}{lrrrr}
    \toprule
    Method & Test Acc & Trainable Params & Peak Mem (MB) & Time (s) \\
    \midrule
    Linear probe              & 37.95\% &     17,325 &   260.7 & 105.4 \\
    LoRA ($r=8$, $\alpha=16$) & 41.44\% &    275,373 & 1,210.7 & 139.0 \\
    Full fine-tuning          & 63.78\% & 10,755,117 & 1,760.5 & 121.6 \\
    \bottomrule
    \end{tabular}
    \end{center}

    Full fine-tuning achieves the highest accuracy (63.8\%) by a large margin because EuroSAT and RESISC45 differ substantially in class structure (10 vs.\ 45 classes, multispectral satellite vs.\ aerial RGB imagery), so the frozen EuroSAT representation transfers imperfectly and requires intermediate features---not just the final head---to be updated. LoRA outperforms the linear probe by 3.5 percentage points despite training only $16\times$ more parameters (275K vs.\ 17K), showing that adapting the attention Q/V projections yields a more flexible adjustment of the representation than a frozen linear readout. Peak memory follows the amount of stored activations during backpropagation: linear probe requires no gradient to pass through the ViT (261~MB), while LoRA and full fine-tuning both backpropagate through every block (1,211~MB and 1,761~MB), with full fine-tuning's extra overhead coming from larger optimizer states for its 10M trainable parameters. Training times are similar across methods (105--139~s) because all three run the same forward pass through the ViT; the LoRA backward is nearly as expensive as full fine-tuning since gradients must still flow through all blocks to reach the A and B matrices.

    \item lora\_rank:

    \begin{center}
    \includegraphics[width=0.7\textwidth]{../runs/resisc_rank_sweep/rank_sweep.png}
    \end{center}

    Accuracy increases monotonically across all seven ranks (32.7\%, 35.3\%, 38.0\%, 41.1\%, 44.4\%, 47.1\%, 49.5\% for $r = 1, 2, 4, 8, 16, 32, 64$), with per-doubling gains of roughly 2.4--3.3 percentage points. The curve is approximately log-linear. In large language models ($d_\text{model} \gg 10^3$), $r = 8$ or $r = 16$ is sufficient because the fine-tuning update is empirically very low-rank relative to the enormous weight matrices---the few most salient adaptation directions capture almost all of the needed shift. Our ViT has $d_\text{model}=384$ and is adapting across a substantial domain gap (10-class EuroSAT satellite imagery $\to$ 45-class RESISC45 aerial imagery), so the effective rank of the update is higher relative to $d_\text{model}$: the model must adjust a broader subspace of its attention projections to separate 45 visually diverse classes, and even $r=64$ (one-sixth of $d_\text{model}$) has not exhausted the useful adaptation directions.

\end{enumerate}
\end{question}

\begin{question}{Vision-Language Model}
\begin{enumerate}[label=(\alph*)]

    \item projector: A single linear layer is a fixed affine map whose capacity is entirely spent on rotating and scaling the image-feature subspace into decoder space; it cannot learn curved boundaries between visual-feature clusters, so semantically distinct image regions that land in different parts of the ViT's representation space may still be indistinguishable after projection. A hidden GELU layer introduces nonlinear capacity precisely at the modality boundary, which is critical when both the ViT encoder and the LM decoder are frozen: the projector must single-handedly bridge two independently pretrained representation spaces, and a shallow nonlinearity lets it learn which dimensions of the visual embedding correspond to decoder concepts without altering either frozen model.

    \item injection\_compare: We train three configurations for 2000 steps each (freeze~A: projector only, causal mask), reporting the best validation accuracy on 500 held-out CLEVR examples.

    \begin{center}
    \begin{tabular}{lrrrr}
    \toprule
    Strategy & Val Acc & Vis.\ Tokens & Peak Mem (MB) & Time/step (s) \\
    \midrule
    CLS-only prefix     & 26.2\% & 1  & 7\,922.5  & 0.63 \\
    All-patches prefix  & 14.8\%$^\dagger$ & 65 & 11\,340.9 & 0.74 \\
    Interleaved         &  1.8\% & 65 & 11\,284.0 & 0.74 \\
    \bottomrule
    \multicolumn{5}{l}{\footnotesize $\dagger$~500-step result (overwritten by masking cell; full 2000-step run not recovered).}
    \end{tabular}
    \end{center}

    The CLS strategy achieves the highest accuracy (26.2\%) despite providing the least visual information. With only one visual token, the projector needs only to map a single global summary vector into decoder space — a relatively easy one-dimensional problem that 2000 steps is sufficient to solve. The all-patches prefix provides 65 tokens (64 patches $+$ CLS for a $64\!\times\!64$ image with $P\!=\!8$), requiring the projector to learn a useful mapping for 65 spatial positions simultaneously; at 500 steps (cut short by a naming collision with the masking cell), it already reaches 14.8\%, suggesting it would surpass CLS given more training. The interleaved mode persistently fails (1.8\% over 2000 steps): placing 65 visual tokens immediately after BOS violates the frozen decoder's strong prior that BOS is followed by ordinary text, producing chaotic internal representations for the question tokens and preventing coherent generation. These results mirror the \texttt{vit\_pooling} finding — CLS is the weakest representation for spatial tasks in principle, but it offers the simplest optimization target for the projector in a short-training regime, while the patch-based strategies need more steps or decoder adaptation to realise their theoretical advantage.

    \item masking: The two attention masks for a sequence of 4 visual tokens ($v_1$--$v_4$) followed by 3 text tokens ($t_1$--$t_3$) are shown below (shaded cells = allowed positions).

    \textbf{(M1) Fully causal:}

    \begin{center}
    \setlength{\tabcolsep}{4pt}
    \begin{tabular}{c|ccccccc}
     & $v_1$ & $v_2$ & $v_3$ & $v_4$ & $t_1$ & $t_2$ & $t_3$ \\\hline
    $v_1$ & \cellcolor{gray!35}1 & 0 & 0 & 0 & 0 & 0 & 0 \\
    $v_2$ & \cellcolor{gray!35}1 & \cellcolor{gray!35}1 & 0 & 0 & 0 & 0 & 0 \\
    $v_3$ & \cellcolor{gray!35}1 & \cellcolor{gray!35}1 & \cellcolor{gray!35}1 & 0 & 0 & 0 & 0 \\
    $v_4$ & \cellcolor{gray!35}1 & \cellcolor{gray!35}1 & \cellcolor{gray!35}1 & \cellcolor{gray!35}1 & 0 & 0 & 0 \\
    $t_1$ & \cellcolor{gray!35}1 & \cellcolor{gray!35}1 & \cellcolor{gray!35}1 & \cellcolor{gray!35}1 & \cellcolor{gray!35}1 & 0 & 0 \\
    $t_2$ & \cellcolor{gray!35}1 & \cellcolor{gray!35}1 & \cellcolor{gray!35}1 & \cellcolor{gray!35}1 & \cellcolor{gray!35}1 & \cellcolor{gray!35}1 & 0 \\
    $t_3$ & \cellcolor{gray!35}1 & \cellcolor{gray!35}1 & \cellcolor{gray!35}1 & \cellcolor{gray!35}1 & \cellcolor{gray!35}1 & \cellcolor{gray!35}1 & \cellcolor{gray!35}1 \\
    \end{tabular}
    \end{center}

    \textbf{(M2) Bidirectional image, causal text:}

    \begin{center}
    \setlength{\tabcolsep}{4pt}
    \begin{tabular}{c|ccccccc}
     & $v_1$ & $v_2$ & $v_3$ & $v_4$ & $t_1$ & $t_2$ & $t_3$ \\\hline
    $v_1$ & \cellcolor{gray!35}1 & \cellcolor{gray!35}1 & \cellcolor{gray!35}1 & \cellcolor{gray!35}1 & 0 & 0 & 0 \\
    $v_2$ & \cellcolor{gray!35}1 & \cellcolor{gray!35}1 & \cellcolor{gray!35}1 & \cellcolor{gray!35}1 & 0 & 0 & 0 \\
    $v_3$ & \cellcolor{gray!35}1 & \cellcolor{gray!35}1 & \cellcolor{gray!35}1 & \cellcolor{gray!35}1 & 0 & 0 & 0 \\
    $v_4$ & \cellcolor{gray!35}1 & \cellcolor{gray!35}1 & \cellcolor{gray!35}1 & \cellcolor{gray!35}1 & 0 & 0 & 0 \\
    $t_1$ & \cellcolor{gray!35}1 & \cellcolor{gray!35}1 & \cellcolor{gray!35}1 & \cellcolor{gray!35}1 & \cellcolor{gray!35}1 & 0 & 0 \\
    $t_2$ & \cellcolor{gray!35}1 & \cellcolor{gray!35}1 & \cellcolor{gray!35}1 & \cellcolor{gray!35}1 & \cellcolor{gray!35}1 & \cellcolor{gray!35}1 & 0 \\
    $t_3$ & \cellcolor{gray!35}1 & \cellcolor{gray!35}1 & \cellcolor{gray!35}1 & \cellcolor{gray!35}1 & \cellcolor{gray!35}1 & \cellcolor{gray!35}1 & \cellcolor{gray!35}1 \\
    \end{tabular}
    \end{center}

    M2 is expected to perform better because it preserves the ViT's original bidirectional view of the image: every visual token can aggregate information from all other visual tokens, matching the encoder's design. M1 imposes an arbitrary rasterisation order on 2-D patches, causing earlier patches to have incomplete context and later patches to be influenced by irrelevant earlier patches in the sequence. Additionally, M2 prevents visual tokens from attending to text tokens (which they haven't ``seen'' yet), matching the causal boundary at the image-to-text interface.

    \begin{center}
    \begin{tabular}{lrr}
    \toprule
    Mask mode & Val Acc (500 steps) & Time/step (s) \\
    \midrule
    Causal (M1)        & 14.8\% & 0.74 \\
    Image-bidir (M2)   &  1.4\% & 2.05 \\
    \bottomrule
    \end{tabular}
    \end{center}

    Contrary to the theoretical prediction, causal masking (M1) substantially outperforms image-bidir (M2) at 500 steps in the projector-only setting. With a randomly initialised projector, bidirectional visual self-attention in the decoder mixes 65 uninformative patch embeddings with each other, producing chaotic visual representations; the frozen decoder's attention heads — trained exclusively on causal text sequences — amplify this chaos rather than filtering it. The causal mask is closer to what the frozen decoder was pretrained on, so the question-token representations degrade less. M2 is also 2.77$\times$ slower per step due to SDPA with the custom 4D additive mask. Under a freeze configuration that adapts the decoder (B/C/D), M2's theoretical advantage would likely emerge.

    \item freezing: Using all-patches injection with image-bidir masking, we train four configurations for 1500 steps each.

    \begin{center}
    \begin{tabular}{llrrrr}
    \toprule
    Config & Decoder & Val Acc & Trainable & Peak Mem (MB) & Time/step (s) \\
    \midrule
    A: projector only   & frozen      &  1.6\% &      2,066,880 & 11\,340.2 & 2.05 \\
    B: proj.\ + LoRA    & LoRA $r=8$  & 21.0\% &      2,886,080 & 11\,740.8 & 2.13 \\
    C: proj.\ + full FT & full FT     & 28.2\% & 363,888,000    & 16\,132.8 & 2.18 \\
    D: all full FT      & full FT     & 20.8\% & 374,625,792    & 16\,635.2 & 2.25 \\
    \bottomrule
    \end{tabular}
    \end{center}

    Config~C (frozen ViT, full decoder fine-tuning) achieves the best accuracy (28.2\%), confirming that adapting the decoder is the critical bottleneck when the projector-only regime fails. The jump from A to B (1.6\%~$\to$~21.0\%) shows that even sparse LoRA adapters (819,200 extra trainable parameters, 0.22\% of the model) are sufficient to teach the decoder how to process visual-token prefixes, enabling the all-patches representation to realise most of its spatial potential. Full decoder fine-tuning (C, 28.2\%) further improves over LoRA by adjusting every weight in the 363M-parameter SmolLM2 decoder, allowing deeper restructuring of the decoder's internal representations. Config~D (all three components fine-tuned) falls back to 20.8\%: jointly fine-tuning the EuroSAT-pretrained ViT alongside the decoder appears to destabilise the visual encoder within 1500 steps, degrading the patch features that the decoder has just learned to use. Memory scales predictably: frozen decoder (A, B: $\approx$\,11.7~GB) vs.\ full decoder gradients (C, D: $\approx$\,16.4~GB). In the two-stage VLM recipe, A corresponds to Stage~1 (align the projector; cheap) and C to Stage~2 (instruction-tune the decoder; expensive but necessary for strong task performance); B offers a practical middle ground.

    \item vlm\_qualitative: We evaluate the best checkpoint (Config~C, all-patches + image-bidir, 28.2\% overall) on 10 held-out CLEVR examples. Per-type breakdown: exist 56.7\%, compare\_attr 42.9\%, spatial 31.2\%, query\_attr 16.2\%, count 12.9\%.

    \begin{center}
    \begin{minipage}[t]{0.47\textwidth}
      \centering
      \includegraphics[width=\textwidth]{../runs/vlm_qualitative/CLEVR_val_000041.png}\\[2pt]
      {\footnotesize \textit{Does the big object behind the cyan matte cylinder have the same material as the big gray object?}\\[1pt]
      [spatial]\enspace Gold: \textbf{no}\enspace Pred: no\enspace\checkmark}
    \end{minipage}
    \hspace{0.04\textwidth}
    \begin{minipage}[t]{0.47\textwidth}
      \centering
      \includegraphics[width=\textwidth]{../runs/vlm_qualitative/CLEVR_val_000196.png}\\[2pt]
      {\footnotesize \textit{The rubber object that is the same color as the tiny metal ball is what size?}\\[1pt]
      [query\_attr]\enspace Gold: \textbf{small}\enspace Pred: small\enspace\checkmark}
    \end{minipage}

    \medskip

    \begin{minipage}[t]{0.47\textwidth}
      \centering
      \includegraphics[width=\textwidth]{../runs/vlm_qualitative/CLEVR_val_000222.png}\\[2pt]
      {\footnotesize \textit{Is the number of small rubber spheres behind the small yellow object greater than the number of green things behind the big green cylinder?}\\[1pt]
      [spatial]\enspace Gold: \textbf{no}\enspace Pred: no\enspace\checkmark}
    \end{minipage}
    \hspace{0.04\textwidth}
    \begin{minipage}[t]{0.47\textwidth}
      \centering
      \includegraphics[width=\textwidth]{../runs/vlm_qualitative/CLEVR_val_000241.png}\\[2pt]
      {\footnotesize \textit{Are there fewer small purple shiny blocks than cyan matte things?}\\[1pt]
      [count]\enspace Gold: \textbf{no}\enspace Pred: no\enspace\checkmark}
    \end{minipage}

    \medskip

    \begin{minipage}[t]{0.47\textwidth}
      \centering
      \includegraphics[width=\textwidth]{../runs/vlm_qualitative/CLEVR_val_000324.png}\\[2pt]
      {\footnotesize \textit{What number of objects are either big shiny cylinders left of the small yellow metallic block or cyan things right of the small sphere?}\\[1pt]
      [spatial]\enspace Gold: \textbf{3}\enspace Pred: 3\enspace\checkmark}
    \end{minipage}
    \hspace{0.04\textwidth}
    \begin{minipage}[t]{0.47\textwidth}
      \centering
      \includegraphics[width=\textwidth]{../runs/vlm_qualitative/CLEVR_val_000055.png}\\[2pt]
      {\footnotesize \textit{How big is the green metal object right of the cube that is right of the tiny thing to the left of the cyan cylinder?}\\[1pt]
      [spatial]\enspace Gold: \textbf{small}\enspace Pred: 3\enspace\texttimes}
    \end{minipage}

    \medskip

    \begin{minipage}[t]{0.47\textwidth}
      \centering
      \includegraphics[width=\textwidth]{../runs/vlm_qualitative/CLEVR_val_000149.png}\\[2pt]
      {\footnotesize \textit{There is a red thing; does it have the same shape as the purple shiny object behind the big green rubber object?}\\[1pt]
      [spatial]\enspace Gold: \textbf{yes}\enspace Pred: no\enspace\texttimes}
    \end{minipage}
    \hspace{0.04\textwidth}
    \begin{minipage}[t]{0.47\textwidth}
      \centering
      \includegraphics[width=\textwidth]{../runs/vlm_qualitative/CLEVR_val_000259.png}\\[2pt]
      {\footnotesize \textit{Is the thing behind the small blue object made of the same material as the object to the left of the cube?}\\[1pt]
      [spatial]\enspace Gold: \textbf{yes}\enspace Pred: no\enspace\texttimes}
    \end{minipage}

    \medskip

    \begin{minipage}[t]{0.47\textwidth}
      \centering
      \includegraphics[width=\textwidth]{../runs/vlm_qualitative/CLEVR_val_000275.png}\\[2pt]
      {\footnotesize \textit{How many cylinders are gray metallic objects or red things?}\\[1pt]
      [count]\enspace Gold: \textbf{2}\enspace Pred: \texttt{: 0}\enspace\texttimes}
    \end{minipage}
    \hspace{0.04\textwidth}
    \begin{minipage}[t]{0.47\textwidth}
      \centering
      \includegraphics[width=\textwidth]{../runs/vlm_qualitative/CLEVR_val_000284.png}\\[2pt]
      {\footnotesize \textit{There is a small object in front of the small object behind the matte object in front of the tiny blue matte thing; what shape is it?}\\[1pt]
      [spatial]\enspace Gold: \textbf{cylinder}\enspace Pred: sphere\enspace\texttimes}
    \end{minipage}
    \end{center}

    The model is strongest on existence and binary attribute-comparison questions (56.7\% and 42.9\%), where a yes/no response can often be derived from coarse scene statistics without precise spatial localisation; it is weakest on counting (12.9\%) and attribute queries (16.2\%), which require identifying specific objects and their properties in a scene the encoder was never pretrained on. Among the five incorrect cases, four (examples 7, 8, 10, and the shape-prediction part of example 6) are plausibly encoder failures: the EuroSAT-pretrained ViT produces patch features calibrated to satellite-image textures and land-use patterns, not to the material, shape, and colour of synthetic 3D objects, so the decoder receives no useful signal to distinguish ``yes'' from ``no'' or ``cylinder'' from ``sphere.'' Example~9 is a decoder failure of a different kind: the prediction ``\texttt{: 0}'' includes a stray colon copied from the ``\texttt{Answer:}'' suffix, suggesting that the answer-only label masking occasionally leaks prompt structure into generation. To disentangle the two failure modes experimentally, one could (a) replace the image with a ground-truth text scene description and measure decoder-only accuracy --- if it is near perfect, the encoder is the bottleneck --- or (b) swap the ViT for a CLIP ViT-L pretrained on natural images and measure the accuracy gain; a large gain implicates the encoder, while a small gain implicates the decoder's limited instruction-following with frozen or lightly adapted weights.

\end{enumerate}
\end{question}

% ===========================================================================
\begin{question}{6 Positional Encodings and RoPE}
\begin{enumerate}[label=\textbf{\arabic*.}]

    \item rope\_1d: RoPE is a rotation applied independently to each
    $(x_{2i},\,x_{2i+1})$ pair; a 2-D rotation is a norm-preserving map,
    so the L2 norm of the full $d$-dimensional vector is the root-sum-of-squares
    of $d/2$ pair norms, each of which is unchanged.
    Empirically, measuring
    $\max_{b,h,t}\bigl|\|x_{b,h,t}\|-\|\mathrm{RoPE}(x)_{b,h,t}\|\bigr|$
    on a random $(2,4,16,64)$ tensor at positions $0\ldots15$ yields a
    maximum absolute deviation of $\approx 3\times10^{-7}$, consistent with
    IEEE 754 float32 round-off and confirming the rotation preserves norms
    to within numerical precision.

    \item rope\_vs\_learned:

    \begin{center}
    \begin{tabular}{lrr}
    \toprule
    Positional encoding & Train-size val acc & Extrap.\ val acc (96$\times$96) \\
    \midrule
    Learned PE          & (pending) & (pending) \\
    1D RoPE             & (pending) & (pending) \\
    \bottomrule
    \end{tabular}
    \end{center}

    Both models are trained CLIP-style on EuroSAT for 20 epochs at $64\times64$
    ($8\times8=64$ patches) and then evaluated at $96\times96$ ($12\times12=144$
    patches).
    For the learned-PE baseline the patch positional embeddings are bilinearly
    interpolated from the $8\times8$ training grid to the $12\times12$
    evaluation grid before being added to the patch tokens; the CLS positional
    embedding is kept unchanged.
    Learned PE is expected to degrade at the extrapolation size because the
    model never sees position indices beyond 64 during training, and bilinear
    interpolation only approximates the learned geometry.
    1D RoPE should extrapolate more gracefully because its rotation angles are
    computed from a fixed formula ($\theta_i = 10000^{-2i/d}$) that is
    well-defined at any position index, and the relative-position property
    ($q_m\cdot k_n$ depends only on $m-n$) is preserved regardless of absolute
    position magnitude.

    \item rope\_2d:

    \begin{center}
    \begin{tabular}{lrr}
    \toprule
    Positional encoding & Train-size val acc & Extrap.\ val acc (96$\times$96) \\
    \midrule
    Learned PE          & (pending) & (pending) \\
    1D RoPE             & (pending) & (pending) \\
    2D RoPE             & (pending) & (pending) \\
    \bottomrule
    \end{tabular}
    \end{center}

    2D RoPE splits the head dimension in half: the first $d/2$ channels rotate
    by the patch's column index and the second $d/2$ channels by the row index,
    so the attention dot product between two patches depends on both their
    horizontal and vertical relative offsets.
    On EuroSAT, which contains spatially structured land-use patterns (rivers
    are horizontal, roads diagonal, etc.), the 2D inductive bias is expected to
    help marginally over 1D RoPE at train size; at the extrapolation size it
    should remain robust since coordinates are still drawn from the fixed
    frequency schedule.

    \item mrope\_written:

    \textbf{(1) What goes wrong with naive 1D position IDs?}
    With naive 1D assignment the 65 visual tokens (CLS at 0, patches at
    1--64) occupy positions 0--64, pushing the 50 text tokens to positions
    65--114.
    (a)~SmolLM2 is trained on text sequences up to 2048 tokens, so positions
    65--114 are well inside its training distribution for short contexts, but
    with larger images (e.g.\ LLaVA's $16\times16=256$ patches) text tokens
    would reach positions 257+, which can exceed practical training lengths and
    degrade the quality of the decoder's rotary representations for those
    positions.
    (b)~More critically, treating the $8\times8$ patch grid as a flat 1D
    sequence destroys its 2D structure: patch~$(0,0)$ at position~1 appears
    ``adjacent'' to patch~$(0,1)$ at position~2, but patch~$(7,7)$ at
    position~64 is also made to appear ``close'' to patch~$(7,6)$ at
    position~63 while being made ``far'' from patch~$(6,7)$ at position~56,
    even though those two patches are equally separated in the image.
    The decoder's RoPE thus cannot distinguish left-right from up-down
    relations, limiting spatial reasoning.

    \textbf{(2) Under M-RoPE, what position does the first text token get?}
    In M-RoPE each token carries a 3D position $(t, h, w)$.
    Image patches at grid cell $(r, c)$ receive $(t=0,\, h=r,\, w=c)$; the
    temporal index is fixed at 0 for all image tokens.
    The first text token following the image block gets
    $(t=G,\, h=G,\, w=G)$, where $G$ is the grid side length (e.g.\ $G=8$
    for a $64\times64$ image with patch size~8).
    This choice is sensible for two reasons: first, it ensures text tokens
    begin at a temporal position strictly greater than any image token's
    temporal coordinate (0), so the model can unambiguously distinguish visual
    from linguistic context via the temporal channel; second, text position IDs
    start at $G$ rather than $G^2+1$ (the naive total), which keeps them
    squarely within the range the decoder was pretrained on and reduces the
    effective positional gap between the image and the first generated word.

    \textbf{(3) Why three chunks rather than two?}
    M-RoPE splits the attention head dimension into three equal chunks, one
    rotated by $t$, one by $h$, and one by $w$.
    If the temporal index $t$ were dropped and only $(h,w)$ were used, text
    tokens would have no natural $(h,w)$ coordinates: they form a 1D sequence,
    so one common assignment is $(h=i,\, w=i)$ for the $i$-th text token.
    But this conflates the text's sequential distance with spatial image
    distance: the relative $(h,w)$ offset between text token $i$ and patch
    $(r,c)$ would be $(i-r,\, i-c)$, a direction that mixes image rows/columns
    with text token indices in a semantically meaningless way.
    The temporal index provides a dedicated channel where text tokens can track
    their sequential order independently of image coordinates, while image tokens
    share a fixed $t=0$ that keeps them distinguishable from all text tokens in
    at least one dimension.

    \item mrope\_impl (bonus):

    \begin{center}
    \begin{tabular}{lrr}
    \toprule
    Position assignment & Overall val acc & Spatial val acc \\
    \midrule
    Naive 1D            & (pending) & (pending) \\
    M-RoPE-style        & (pending) & (pending) \\
    \bottomrule
    \end{tabular}
    \end{center}

    Both models use the best §5 configuration (all-patches injection,
    image-bidir masking, freeze config~C) and are trained for 1500 steps.
    In the M-RoPE variant, the CLS token receives position~0; each patch at
    grid cell $(r,c)$ of an $8\times8$ grid receives position
    $\max(r,c)+1 \in \{1,\ldots,8\}$; and the $k$-th text token receives
    position $9+k$.
    This is implemented by passing custom \texttt{position\_ids} to the
    decoder's \texttt{forward()} call, overriding SmolLM2's default
    $0,1,\ldots,T{-}1$ assignment.
    The key benefit is that text tokens start at position~9 instead of~65,
    while spatially adjacent patches receive more similar position IDs
    (reflecting 2D proximity) than under the naive scheme.
    Whether this compact assignment helps on CLEVR spatial questions will
    depend on whether the model's capacity bottleneck is the positional
    encoding or the EuroSAT-pretrained encoder features.

\end{enumerate}
\end{question}

\end{document}
